\section{Introduction}
\label{sec:intro}

Les correcteurs d'ordre fractionnaire promettent un réglage plus fin que leurs
homologues entiers, au prix de degrés de liberté supplémentaires que l'on
confie volontiers à des métaheuristiques. Pour un convertisseur de puissance,
cette promesse se heurte à trois réalités de mise en œuvre : l'échantillonnage
et le retard, la saturation d'actionneur et l'anti-emballement, et
l'approximation rationnelle des opérateurs fractionnaires.

L'étude part d'un échec documenté. Une campagne FO-LQI $\to$ FO-PIDF sur un
Buck non idéal de 48~V a consommé plus de 1800 évaluations réparties sur six
algorithmes sans jamais atteindre la faisabilité robuste. Le diagnostic
retenu à l'époque mettait en cause l'écrêtage du gain proportionnel projeté
sur la borne $\log_{10}k_p=-4$, atteinte dans $79~\%$ des candidats.

\paragraph{Question de méthode.}
Lorsqu'une recherche paramétrique échoue, comment déterminer si l'échec
vient de l'algorithme, de la paramétrisation, ou du problème lui-même ? Nous
répondons par un protocole en trois temps : un audit de saturation des bornes
qui sépare saturations structurelles et pathologiques ; des sondes conjointes
dans l'espace du correcteur, qui testent l'accessibilité de la région
faisable indépendamment de la paramétrisation ; une campagne confirmatoire
préenregistrée qui mesure l'échec avec les statistiques qui conviennent, au
lieu de le supposer.

\paragraph{Contributions.}
\begin{enumerate}
\item Un audit de saturation des bornes de paramétrisation, qui montre ici
  que l'écrêtage de $k_p$ n'est pas la cause de l'échec antérieur.
\item Une chaîne FO-LQI $\to$ FO-PIDF reconstruite et vérifiée : la
  projection reproduit l'identité exacte en ordre entier et mesure son
  erreur hors de ce domaine.
\item Une campagne confirmatoire préenregistrée en trois bras, dont les
  résultats négatifs sont rapportés intégralement, graines stériles comprises.
\item L'identification d'un verrou physique indépendant du correcteur (le
  plancher de tension à rapport cyclique minimal) et d'un angle mort du
  cahier des charges (les cycles limites invisibles aux contraintes).
\end{enumerate}
